\section*{अध्याय \ref{ch:b2:wilson-cycle} --- प्राचीन पुंजक और विल्सन चक्र}

\begin{solution}{exo:b2:wilson-cycle:1}
(क) $600 - 400 = \qty{200}{Myr}$, किनारों की दरार से अंतिम टक्कर तक।
(ख) $480 - 400 = \qty{80}{Myr}$। (ग) कैलेडोनियाई पर्वतजनन चरणों का एक समुच्चय है जो
कोई 200 मिलियन वर्षों में फैला है, और हर चरण भिन्न स्थानों पर भिन्न समयों पर आरंभ हुआ: यह शब्द
एक महासागर और उसकी पट्टी की स्थिति बताता है, किसी क्षण की नहीं। आयु स्तरिक मापक्रम से या
मिलियन वर्षों में दी जाती है।
\end{solution}

\begin{solution}{exo:b2:wilson-cycle:2}
(ग) पूर्वी अफ़्रीकी दरार, दरार-चरण; (क) लाल सागर, एक युवा महासागर; (घ) अटलांटिक, निष्क्रिय किनारों वाला एक चौड़ा
महासागर; (ङ) दक्षिण अमेरिका का प्रशांत किनारा, क्षेपण और किसी महासागर का
बंद होना; (ख) हिमालय, टक्कर; (च) ब्रिटनी, समतल घिसी एक पट्टी, जिसके बाद
फिर से दरार आरंभ हो सकती है।
\end{solution}

\begin{solution}{exo:b2:wilson-cycle:3}
(क) आरंभिक सिल्यूरियन परतों के जमने के बाद, यानी \qty{443.1}{Ma} के बाद, और
उत्तर डेवोनियन परतों के जमने से पहले, यानी \qty{358.9}{Ma} से पहले, डेवोनियन के अंत से।
(ख) $443.1 - 358.9 = \qty{84.2}{Myr}$ से कम; वास्तव में और कम, क्योंकि निचली परतें
सिल्यूरियन के आधार से नई हैं और ऊपरी परतें उसके शीर्ष से पुरानी। (ग) महासागर में
धूसर गंदलाव-निक्षेपों का जमना; कैलेडोनियाई टक्करों में वलन और विदलन; उत्थान और
अपरदन से एक असमतल स्थल-सतह; लाल संकोणाश्मों और संगुटिकाश्मों का जमना; वैरिस्की पर्वतजनन में मंद
झुकाव; आज का अपरदन, जो खड़ी चट्टान को काटता है।
\end{solution}

\begin{solution}{exo:b2:wilson-cycle:4}
(क) दोनों: महाद्वीपीय भूपर्पटी का एक खंड जो आयपेटस को पार करते समय महासागरों द्वारा महाद्वीपों से अलग था,
और, जुड़ जाने के बाद, भ्रंशों से घिरा एक खंड जिसका इतिहास उसके
पड़ोसियों से भिन्न है। (ख) कोई नहीं: एक महाद्वीप, किसी महाद्वीप से जोड़ा गया खंड नहीं। (ग) \omterm{def:b2:wilson-cycle:terrane}{टेरेन}, \omterm{def:b2:wilson-cycle:terrane}{सूक्ष्म-महाद्वीप}
नहीं: उसकी भूपर्पटी चाप की है, किसी महाद्वीप का टुकड़ा नहीं। (घ) कोई नहीं: यह
\omterm{def:b2:wilson-cycle:terrane}{टेरेनों} के बीच की सीमा है, बंद हुए महासागर का चिह्न।
\end{solution}

\begin{solution}{exo:b2:wilson-cycle:5}
(क) $3000 - 2600 = 400$, $2600 - 1500 = 1100$, $1500 - 875 = 625$, $875 - 600 = 275$, $600 -
262.5 = 337.5$; माध्य $2737.5/5 = \qty{547.5}{Myr}$। (ख) $1100/275 = 4.0$। (ग) $(1500 -
262.5)/3 = \qty{412.5}{Myr}$। संदिग्ध बिंदु हटाने पर माध्य एक-चौथाई खिसक जाता है
और अंतराल चार गुना तक बदलते हैं: एक मोटी पुनरावृत्ति, आवर्तकाल नहीं।
\end{solution}

\begin{solution}{exo:b2:wilson-cycle:6}
(क) $3000/(70 - 25) = \qty{66.7}{Myr}$। (ख) $3000/70 = \qty{42.9}{Myr}$। (ग) कटक के
क्षेपित हो जाने के बाद बचा महासागर-तल दूर की प्लेट का भाग है, जो दूर के
महाद्वीप को ले जाती है: बंद होने की दर खाई के सापेक्ष उस प्लेट का वेग है।
\end{solution}

\begin{solution}{exo:b2:wilson-cycle:7}
(क) $700/8 = \qty{87.5}{km}$; $700/30 = \qty{23.3}{km}$। (ख) वैरिस्की पट्टी:
\qty{30}{\celsius} प्रति किलोमीटर पर मध्य भूपर्पटी गलती है; \qty{8}{\celsius} प्रति किलोमीटर पर
गलन-तापमान मोटी हुई भूपर्पटी के आधार से भी नीचे होता है। (ग) निम्न प्रवणता
उन चट्टानों का वर्णन करती है जो किसी क्षेपण-क्षेत्र में तेज़ी से नीचे ले जाई गईं, इससे पहले कि वे गर्म हो सकें। जो चट्टानें
मोटी हुई भूपर्पटी में गहरी रहती हैं वे लाखों वर्षों में सामान्य जियोथर्म की ओर गर्म होती हैं, और
तब कम गहराई पर गलन तक पहुँच सकती हैं।
\end{solution}

\begin{solution}{exo:b2:wilson-cycle:8}
(क) $\tau\ln(3/0.3) = 26.3\times2.303 = \qty{60.6}{Myr}$। (ख) $290 - 60.6 = \qty{229}{Ma}$,
उत्तर ट्राइऐसिक में। (ग) $\tau = (290 - 251.9)/\ln10 = \qty{16.5}{Myr}$। (घ) मोटी हुई भूपर्पटी का
विस्तारी निपात, जो उसे पतला करता है और उसकी सतह को सीधे नीचे लाता है; साथ ही
अधिक अपरदन-दर, अधिक आर्द्र जलवायु या कोमल चट्टानें, या कम आरंभिक ऊँचाई।
\end{solution}

\begin{solution}{exo:b2:wilson-cycle:9}
(क) सीवन, यानी आयपेटस का चिह्न, ब्रिटेन और आयरलैंड में स्थल पर, अटलांटिक के
पूर्व में, स्थित है; आयपेटस के दूर वाले पक्ष से आई लॉरेंशिया की भूपर्पटी रॉकॉल के पास
यूरोपीय पक्ष पर स्थित है। नई दरार पुराने सीवन के पश्चिम में कटी: वह उसके निकट खुली, उसके अनुदिश नहीं।
(ख) सीवन के उत्तर में लॉरेंशियाई किनारे का आधारशैल और आवरण; उसके दक्षिण में
एवलोनिया के, भिन्न इतिहासों के साथ और, आरंभिक पैलियोज़ोइक में, भिन्न जीवों के साथ। (ग) दरार
80 मिलियन वर्ष बाद आई, जो उन 30 से 50 से अधिक है जिनके बाद कोई ठंडा हुआ मंडल
दृढ़ हो चुका हो सकता है; फिर भी मध्य अटलांटिक ने ऐपलेशियन का अनुसरण किया। पट्टी की ऊष्मा मिट जाती है; उसके
भ्रंश, संरचनाएँ और संघटन बने रहते हैं और अब भी किसी दरार को दिशा दे सकते हैं।
\end{solution}

\begin{solution}{exo:b2:wilson-cycle:10}
(क) $70 - 40/2 = \qty{50}{mm}$ प्रति वर्ष। (ख) $\qty{1500}{km}/\qty{50}{km.Myr^{-1}} =
\qty{30}{Myr}$। (ग) महासागर $70 - 40 = \qty{30}{mm}$ प्रति वर्ष से संकरा होता है: $3000 - 30\times30 =
\qty{2100}{km}$। (घ) तब युवा, गर्म स्थलमंडल खाई पर पहुँचता है, और कोई पुराना, ठंडा स्लैब
चाप के नीचे अब निर्जलित नहीं होता: चाप का मैग्मावाद बदल जाता है या रुक जाता है।
\end{solution}

\begin{solution}{exo:b2:wilson-cycle:11}
(क) $h = \Delta T(\rho_m - \rho_c)/\rho_m = 25\times600/3300 = \qty{4.5}{km}$। (ख) पूरी अतिरिक्त
\qty{25}{km}, पर्वतमाला की ऊँचाई की 5.5 गुनी। (ग) $2700\times9.8\times25\,000 =
\qty{6.6e8}{Pa}$, यानी लगभग \qty{0.66}{GPa}। (घ) लगभग \qty{0.7}{GPa} पर और मोटी हुई भूपर्पटी की
जड़ के तापमानों पर, यानी कई सौ डिग्री पर, एम्फ़ीबोलाइट या ग्रैन्यूलाइट
संलक्षणी: पुराने पुंजकों के नाइस और \omterm{def:b2:crustal-melting:migmatite}{मिग्मेटाइट}।
\end{solution}

\begin{solution}{exo:b2:wilson-cycle:12}
(क) $S = 30\times10 = \qty{300}{km}$; $W_0 = 300 + 300 = \qty{600}{km}$; $T = 35\times600/300 =
\qty{70}{km}$। (ख) $h = 35\times600/3300 = \qty{6.36}{km}$; $r = 4.5h = \qty{28.6}{km}$।
(ग) $\Delta F = 2700\times9.8\times6364\times(35\,000 + 3182 + 14\,318) =
\qty{8.8e12}{N.m^{-1}}$, तिब्बत का 1.4 गुना। (घ) ऐसी पट्टी तिब्बत से भारी है, जो पहले ही
फैल रहा है: \omterm{def:b2:himalaya-andes:collapse}{गुरुत्वीय निपात}, अपने ऊँचे क्षेत्रों के विस्तारण और अपने किनारों पर अधिक्षेपण द्वारा।
\end{solution}

\begin{solution}{pb:b2:wilson-cycle:1}
\textbf{1.} $720/(540 - 500) = \qty{18}{km}$ प्रति मिलियन वर्ष, यानी \qty{18}{mm} प्रति वर्ष।
\textbf{2.} \qty{36}{mm} प्रति वर्ष: एक धीमा प्रसार कटक (20 से \qty{55}{mm} प्रति वर्ष)।
\textbf{3.} $36\times(580 - 540) = \qty{1440}{km}$।
\textbf{4.} $36\times(580 - 500) = \qty{2880}{km}$।
\textbf{5.} \qty{580}{Ma} पर शून्य चौड़ाई से स्थिर दर पर सममित प्रसार, और
\qty{500}{Ma} से पहले कोई क्षेपण नहीं।
\textbf{6.} \qty{80}{Myr}, महासागर के दोनों किनारों पर, महाद्वीपीय किनारों से सटा। यह महासागर का
सबसे ठंडा, सबसे मोटा और सबसे सघन महासागरीय स्थलमंडल है, धँसने को सबसे अधिक तत्पर।
\textbf{7.} $64 - 36 = \qty{28}{mm}$ प्रति वर्ष।
\textbf{8.} $2880/28 = \qty{102.9}{Myr}$, \qty{500}{Ma} के बाद: \qty{397.1}{Ma} पर।
\textbf{9.} कटक, जो खाई से \qty{1440}{km} दूर है, उसके निकट $64 - 18 =
\qty{46}{mm}$ प्रति वर्ष से आता है: $1440/46 = \qty{31.3}{Myr}$ बाद, \qty{468.7}{Ma} पर।
\textbf{10.} $2880 - 28\times31.3 = \qty{2004}{km}$।
\textbf{11.} $2004/40 = \qty{50.1}{Myr}$ बाद: $468.7 - 50.1 = \qty{418.6}{Ma}$ पर।
\textbf{12.} स्कैंडी टक्कर के 425 से \qty{400}{Ma} के भीतर।
\textbf{13.} $40\times8 = \qty{320}{km}$।
\textbf{14.} $W_0 = 400 + 320 = \qty{720}{km}$; $T = 35\times720/400 = \qty{63.0}{km}$।
\textbf{15.} $h = 28\times600/3300 = \qty{5.09}{km}$; $r = 4.5\times5.09 = \qty{22.9}{km}$।
\textbf{16.} $\Delta F = 2700\times9.8\times5091\times(35\,000 + 2545 + 11\,455) =
\qty{6.60e12}{N.m^{-1}}$, तिब्बत की ऊर्जा से लगभग \qty{8}{\percent} अधिक।
\textbf{17.} उठना बंद करना और फैलना: अपने ऊँचे भीतरी भाग में विस्तारण और सामान्य भ्रंश,
अपने किनारों पर अधिक्षेप, जैसा तिब्बत आज करता है।
\textbf{18.} $f = 600/3300 = 0.182$; $k = 2.303/11 = 0.209$ प्रति मिलियन वर्ष; $\tau = 1/(kf)
= \qty{26.3}{Myr}$।
\textbf{19.} $3.0\,e^{-30/26.3} = \qty{0.96}{km}$।
\textbf{20.} $26.3\ln(3.0/0.5) = \qty{47.1}{Myr}$, \qty{410}{Ma} के बाद: \qty{362.9}{Ma} पर।
\textbf{21.} अतिरिक्त भूपर्पटी $h/f$ है: $3.0/0.182 = \qty{16.5}{km}$ घटकर $0.5/0.182 =
\qty{2.75}{km}$ रह जाती है, अतः \qty{13.75}{km} चट्टान हटती है।
\textbf{22.} $2700\times9.8\times13\,750 = \qty{3.6e8}{Pa}$, यानी लगभग \qty{0.36}{GPa}।
\textbf{23.} हाँ: उत्तर डेवोनियन तक प्रतिरूप-पट्टी लगभग \qty{1}{km} माध्य
ऊँचाई तक आ चुकी है और अब भी गिर रही है, अपनी घाटियों और कटकों के बीच बचे उच्चावच के साथ; लाल बजरी
\qty{358.9}{Ma} से बहुत पहले उसके निचले भागों में एकत्र हो सकती है, जैसे सिकर पॉइंट पर, जहाँ वह
उच्चावच वाली सतह पर टिकी है।
\textbf{24.} \textbf{\qty[text-series-to-math=true]{2880}{km}}, \qty{500}{Ma} पर, जब
क्षेपण आरंभ हुआ।
\end{solution}
